%  
% Latex Problem Set Template by Sachin Padmanabhan
% I created this when I was a freshman in CS 103,
% and I continue to use it to this day.
% 
% Hope you enjoy!
% 
% There may be problems with this template.
% If so, feel free to contact me.
% 

% Updated for Summer 2019 by Amy Liu

\documentclass{article}
\usepackage{amsmath}
\usepackage{amssymb}
\usepackage{amsthm}
\usepackage{amssymb}
\usepackage{mathdots}
\usepackage[pdftex]{graphicx}
\usepackage{fancyhdr}
\usepackage[margin=1in]{geometry}
\usepackage{multicol}
\usepackage{bm}
\usepackage{listings}
\PassOptionsToPackage{usenames,dvipsnames}{color}  %% Allow color names
\usepackage[table, xcdraw, usenames,dvipsnames]{xcolor}
\usepackage[table]{xcolor}
\usepackage{pdfpages}
\usepackage{algpseudocode}
\usepackage{tikz}
\usepackage{enumitem}
\usepackage[T1]{fontenc}
\usepackage{inconsolata}
\usepackage{framed}
\usepackage{wasysym}
\usepackage[thinlines]{easytable}
\usepackage{hyperref}
\usepackage{minted}
\usepackage{pgfplots}
\usepackage{gensymb}
\usepackage{mathtools}
\usemintedstyle{perldoc}
\usepackage{tkz-euclide}
\usepackage{graphicx}
\usepackage{hyperref}
\usepackage{float}
\usepackage{caption}
\usepackage[most]{tcolorbox}
\usepackage{lipsum}
\usepackage{systeme}

\usetikzlibrary{calc, math, shapes.geometric, shapes.misc, arrows.meta, positioning}

% \usepgfplotslibrary{external}
% \tikzexternalize
\usemintedstyle{perldoc}
\hypersetup{
  colorlinks=true,
  linkcolor=blue,
  filecolor=magenta,
  urlcolor=blue,
}




\title{A First Course in Probability\\Problems}
\author{Leonard Mohr}
\date{August 30, 2024}

\lhead{}
\chead{A First Course In Probability}
\rhead{}
\lfoot{}
\cfoot{My Solutions}
\rfoot{\thepage}

\newcommand{\abs}[1]{\lvert #1 \rvert}
\newcommand{\absfit}[1]{\left\lvert #1 \right\rvert}
\newcommand{\norm}[1]{\left\lVert #1 \right\rVert}
\newcommand{\eval}[3]{\left[#1\right]_{#2}^{#3}}
\renewcommand{(}{\left(}
  \renewcommand{)}{\right)}
\newcommand{\floor}[1]{\left\lfloor#1\right\rfloor}
\newcommand{\ceil}[1]{\left\lceil#1\right\rceil}
\newcommand{\pd}[1]{\frac{\partial}{\partial #1}}
\newcommand{\inner}[1]{\langle#1\rangle}
\newcommand{\cond}{\bigg|}
\newcommand{\rank}[1]{\mathbf{rank}(#1)}
\newcommand{\range}[1]{\mathbf{range}(#1)}
\newcommand{\nullsp}[1]{\mathbf{null}(#1)}
\newcommand{\repr}[1]{\left\langle#1\right\rangle}

\DeclareMathOperator{\Var}{Var}
\DeclareMathOperator{\tr}{tr}
\DeclareMathOperator{\Tr}{\mathbf{Tr}}
\DeclareMathOperator{\diag}{\mathbf{diag}}
\DeclareMathOperator{\dist}{\mathbf{dist}}
\DeclareMathOperator{\prob}{\mathbf{prob}}
\DeclareMathOperator{\dom}{\mathbf{dom}}
\DeclareMathOperator{\E}{\mathbf{E}}
\DeclareMathOperator{\R}{\mathbb{R}}
\DeclareMathOperator{\var}{\mathbf{var}}
\DeclareMathOperator{\quartile}{\mathbf{quartile}}
\DeclareMathOperator{\conv}{\mathbf{conv}}
\DeclareMathOperator{\VC}{VC}
\DeclareMathOperator*{\argmax}{arg\,max}
\DeclareMathOperator*{\argmin}{arg\,min}
\DeclareMathOperator{\Ber}{Bernoulli}
\DeclareMathOperator{\NP}{\mathbf{NP}}
\DeclareMathOperator{\coNP}{\mathbf{coNP}}
\DeclareMathOperator{\TIME}{\mathsf{TIME}}
\DeclareMathOperator{\polytime}{\mathbf{P}}
\DeclareMathOperator{\PH}{\mathbf{PH}}
\DeclareMathOperator{\SIZE}{\mathbf{SIZE}}
\DeclareMathOperator{\ATIME}{\mathbf{ATIME}}
\DeclareMathOperator{\SPACE}{\mathbf{SPACE}}
\DeclareMathOperator{\ASPACE}{\mathbf{ASPACE}}
\DeclareMathOperator{\NSPACE}{\mathbf{NSPACE}}
\DeclareMathOperator{\Z}{\mathbb{Z}}
\DeclareMathOperator{\N}{\mathbb{N}}
\DeclareMathOperator{\EXP}{\mathbf{EXP}}
\DeclareMathOperator{\NEXP}{\mathbf{NEXP}}
\DeclareMathOperator{\NTIME}{\mathbf{NTIME}}
\DeclareMathOperator{\DTIME}{\mathbf{DTIME}}
\DeclareMathOperator{\poly}{poly}
\DeclareMathOperator{\BPP}{\mathbf{BPP}}
\DeclareMathOperator{\ZPP}{\mathbf{ZPP}}
\DeclareMathOperator{\RP}{\mathbf{RP}}
\DeclareMathOperator{\coRP}{\mathbf{coRP}}
\DeclareMathOperator{\BPL}{\mathbf{BPL}}
\DeclareMathOperator{\IP}{\mathbf{IP}}
\DeclareMathOperator{\PSPACE}{\mathbf{PSPACE}}
\DeclareMathOperator{\NPSPACE}{\mathbf{NPSPACE}}
\DeclareMathOperator{\SAT}{\mathsf{SAT}}
\DeclareMathOperator{\NL}{\mathbf{NL}}
\DeclareMathOperator{\PCP}{\mathbf{PCP}}
\DeclareMathOperator{\PP}{\mathbf{PP}}
\DeclareMathOperator{\cost}{cost}
\let\Pr\relax
\DeclareMathOperator*{\Pr}{\mathbf{Pr}}

\makeatletter
\renewcommand*\env@matrix[1][*\c@MaxMatrixCols c]{%
  \hskip -\arraycolsep
  \let\@ifnextchar\new@ifnextchar
  \array{#1}}
\makeatother

% start MZ
\let\oldemptyset\emptyset
\renewcommand{\emptyset}{\text{\O}}
\renewcommand{\labelitemii}{$\bullet$}
\renewcommand\qedsymbol{$\blacksquare$}
\newenvironment{prf}{{\bfseries Proof.}}{\qedsymbol}
\renewcommand{\emph}[1]{\textit{\textbf{#1}}}
\newcommand{\annotate}[1]{\textit{\textcolor{blue}{#1}}}
\usepackage{mdframed}
\usepackage{float}

% For Polyominoes
% Source: https://tex.stackexchange.com/a/110669
\usepackage{color}

\definecolor{lightgray}{rgb}{0.75, 0.75, 0.75}

\newdimen\omsq	\omsq=16pt
\newdimen\omrule    \omrule=1pt
\newdimen\omint

\newif\ifvth	\newif\ifhth	\newif\ifomblank
\def\OMINO#1{%
  \vthtrue \hthtrue
  \vbox{ \offinterlineskip\parindent=0pt \OM#1\relax\vskip1pt}}
\definecolor{gray}{rgb}{0.5, 0.5, 0.5}

\def\OM#1{%
  \omint=\omsq	  \advance\omint-\omrule
  \ifx\relax#1%
  \else
  \ifx\#1 \newline\null \hthtrue \ifvth\vthfalse\else\vskip-\omrule\vthtrue\fi
  \else%
  \ifx .#1\hskip\ifhth \omrule\else \omint\fi
  \else%
  \ifx +#1\def\colour{black}\fi%
  \ifx -#1\def\colour{black}\fi%
  \ifx |#1\def\colour{black}\fi%
  \ifx @#1\def\colour{black}\fi%
  \ifx r#1\def\colour{red}\fi%
  \ifx g#1\def\colour{green}\fi%
  \ifx b#1\def\colour{blue}\fi%
  \ifx y#1\def\colour{yellow}\fi%
  \ifx m#1\def\colour{magenta}\fi%
  \ifx c#1\def\colour{cyan}\fi%
  \ifx x#1\def\colour{lightgray}\fi%
  \textcolor{\colour}{\rule{\ifhth\omrule\else\omsq\fi}{\ifvth\omrule\else\omsq\fi}}%
  \ifhth\else\hskip -\omrule\fi%
  \fi%
  \ifhth\hthfalse\else\hthtrue\fi%
  \fi%
  \expandafter\OM%
  \fi}

\makeatother
% end MZ

\definecolor{shadecolor}{gray}{0.95}

\theoremstyle{plain}
\newtheorem*{lem}{Lemma}

\theoremstyle{plain}
\newtheorem*{claim}{Claim}

\theoremstyle{definition}
\newtheorem*{answer}{Answer}
\newtheoremstyle{indented}
  {1em} % Space above
  {1.5em} % Space below
  {\addtolength{\leftskip}{2.5em}} % Body font
  {} % Indent amount
  {\bfseries} % Theorem head font
  {.} % Punctuation after theorem head
  {.5em} % Space after theorem head
  {} % Theorem head spec (can be left empty, meaning `normal')
%\newtheoremstyle{indented}{\topsep}{\topsep}{\addtolength{\leftskip}{2.5em}}{}{\bfseries}{.}{.5em}{}
\theoremstyle{indented}
\newtheorem{theorem}{Theorem}[section]
\newtheorem{defenition}{Defenition}
\newtheorem*{thm}{Theorem}
\newtheorem{corollary}{Corollary}[theorem]
\newtheorem{lemma}[theorem]{Lemma}

\renewcommand{\headrulewidth}{0.4pt}
\renewcommand{\footrulewidth}{0.4pt}

\setlength{\parindent}{0pt}

\pagestyle{fancy}

\renewcommand{\thefootnote}{\fnsymbol{footnote}}

\makeatletter\@addtoreset{section}{part}\makeatother%

\begin{document}

\maketitle
\tableofcontents

\section{Chapter 1}
\subsection{Problems}\label{subsec:}
\begin{enumerate}
\item
  \begin{enumerate}[label=\alph*)]
\item How many different 7-place license plates are possible if the first 2 places are for letters and the other 5 for numbers?

  \begin{tcolorbox}[%
    enhanced, 
    breakable,
    frame hidden,
    overlay broken = {
      (frame.north west) rectangle (frame.south east);},
    ]
    There are $26$ different letters, and $10$ different numbers.  Since there are five numbers chosen and 2 letters chosen then there are a total of
    \[26^2 \cdot 10^5 = 67,600,000\]
    different possibilities.
  \end{tcolorbox}

  \item Repeat part (a) under the assumption that no letter or number can be repeated in a single license plate. 

    \begin{tcolorbox}[%
      enhanced, 
      breakable,
      frame hidden,
      overlay broken = {
        (frame.north west) rectangle (frame.south east);},
      ]
      The number of possibilities are
      \begin{align*}
26 \cdot 25 \cdot \frac{10!}{5!} &= 19,656,000.
      \end{align*}
    \end{tcolorbox}
  \end{enumerate}

  \item How many outcome sequences are possible when a die is rolled four times, where we say, for instance, that the outcome is $3,4,3,1$ if the first roll landed on $3$, the second on $4$, the third on $3$, and the fourth on $1$?

    \begin{tcolorbox}[%
      enhanced, 
      breakable,
      frame hidden,
      overlay broken = {
        (frame.north west) rectangle (frame.south east);},
      ]
      For a $6$ sided die, there are $6^4$ different possibilities.
    \end{tcolorbox}

    \item Twenty workers are to be assigned to twenty different jobs, one to each job. How many different assignments are possible?

      \begin{tcolorbox}[%
        enhanced, 
        breakable,
        frame hidden,
        overlay broken = {
          (frame.north west) rectangle (frame.south east);},
        ]
        There are $20!$ different possible assignments.
      \end{tcolorbox}

\item John, Jim, Jay, and Jack have formed a band consisting of 4 instruments.  If each of the boys can play all 4 instruments, how many different arrangements are possible? What if John and Jim can play all 4 instruments, but Jay and Jack can each play only piano and drums?

  \begin{tcolorbox}[%
    enhanced, 
    breakable,
    frame hidden,
    overlay broken = {
      (frame.north west) rectangle (frame.south east);},
    ]
    If each member can play all 4 instruments, then there are $4!$ different possible arrangements.  Since Jay and Jack can only plan 2 instruments, then there are $2 \cdot 2 = 4$ different possible arrangements.  
  \end{tcolorbox}

  \item For years, telepthone area codes in the United States and Canada consisted of a sequence of three digits. The first digin was an integer between $2$ and $9$, the second digit was either $0$ or $1$, and the third digit was an integer from $1$ to $9$.  How many area codes were possible? How many area codes starting with a 4 were possible?

    \begin{tcolorbox}[%
      enhanced, 
      breakable,
      frame hidden,
      overlay broken = {
        (frame.north west) rectangle (frame.south east);},
      ]
      There were a total of
      \[ 8 \cdot 2 \cdot 9 = 144\]
      different possible area codes.
      If the first digit was a 4, then there were
      \[1 \cdot 2 \cdot 9 = 18\]
      different possible area codes.
    \end{tcolorbox}

  \item A well-known nursery rhyme starts as follows:
    ``As I was going to St.Ives\\
    I met a man with $7$ wives.\\
    Each wife had $7$ sacks.\\
    Each sack had $7$ cats.\\
    Each cat had $7$ kittens...''
    How many kittens did the traveler meet?

   \begin{tcolorbox}[%
     enhanced, 
     breakable,
     frame hidden,
     overlay broken = {
       (frame.north west) rectangle (frame.south east);},
     ]
     Who knows, the traveler just met the man who had a lot of wifes, not necesarilly the wives and their exorbitant number of sacks containing cats.  But he met,
     \[7^4 = 2401\]
     kittens.
   \end{tcolorbox}

   \item \begin{enumerate}[label=\alph*)]
\item In how many ways can 3 boys and 3 girls sit in a row?

  \begin{tcolorbox}[%
    enhanced, 
    breakable,
    frame hidden,
    overlay broken = {
      (frame.north west) rectangle (frame.south east);},
    ]
    Let's assume that each boy and girl is not unique, all we care about is gender.  First, there are $6!$ possible different arrangements of the boys and girls, but since they are unique, we need to get rid of the $3!$ different combinations of the girls, and the $3!$ different combinations of the boys.  Thus there are
    \[\frac{6!}{3!\cdot 3!}\]
    different possible arrangements.\\\\
    If they are unique, then there are $6! = 720$ different possible arrangements.
  \end{tcolorbox}

\item In how many ways can $3$ boys and $3$ girls sit in a row if the boys and the girls are to sit together?

  \begin{tcolorbox}[%
    enhanced, 
    breakable,
    frame hidden,
    overlay broken = {
      (frame.north west) rectangle (frame.south east);},
    ]
    There are $3!$ ways to arrange boys, and $3!$ ways to arrange the girls.  We also need to select who sits where, thus there are a total of $3! \cdot 3! \cdot 2$ different arrangements.
  \end{tcolorbox}

\item In how many ways if only the boys must sit together?

    \begin{tcolorbox}[%
      enhanced, 
      breakable,
      frame hidden,
      overlay broken = {
        (frame.north west) rectangle (frame.south east);},
      ]
      There are $3!$ ways to arrange boys, and girls, and there are $4$ different places we can place the boys.  Thus there are $3!\cdot 3! \cdot 4$ different possible arrangements.
    \end{tcolorbox}

    \item In how many ways if no two people of the same sex are allowed to sit next to eachother?

      \begin{tcolorbox}[%
        enhanced, 
        breakable,
        frame hidden,
        overlay broken = {
          (frame.north west) rectangle (frame.south east);},
        ]
        This is just alternating boys and girls.  We first choose if a boy or a girl should sit first, then we choose the boys and girls.  Thus we have
        \[2 \cdot 3! \cdot 3! = 72\]
        different options.
      \end{tcolorbox}
    \end{enumerate}

  \item How many different letter arrangements can be made from the letters
    \begin{enumerate}[label=\alph*)]
\item Fluke?

  \begin{tcolorbox}[%
    enhanced, 
    breakable,
    frame hidden,
    overlay broken = {
      (frame.north west) rectangle (frame.south east);},
    ]
    There are $5$ distinct letters, thus $5!$.
  \end{tcolorbox}

  \item Propose?

    \begin{tcolorbox}[%
      enhanced, 
      breakable,
      frame hidden,
      overlay broken = {
        (frame.north west) rectangle (frame.south east);},
      ]
      There are $7$ letters, so there are $7!$ different arrangements of letters, but p, and o are repeated twice each and so we have
      \begin{align*}
\frac{7!}{2!\cdot 2!}
      \end{align*}
      different possibilities.
    \end{tcolorbox}

    \item Mississippi?

      \begin{tcolorbox}[%
        enhanced, 
        breakable,
        frame hidden,
        overlay broken = {
          (frame.north west) rectangle (frame.south east);},
        ]
        There are $11$ characters, but both \texttt{s} and \texttt{i} are repeated $4$ times, and \texttt{p} is repeated two times, therefore
        \[\frac{11!}{4!\cdot 4!\cdot 2!}\]
      \end{tcolorbox}

      \item Arrange?

        \begin{tcolorbox}[%
          enhanced, 
          breakable,
          frame hidden,
          overlay broken = {
            (frame.north west) rectangle (frame.south east);},
          ]
          There are $7$ letters, but both \texttt{r} and \texttt{a} are repeated $2$ times. 
          \[\frac{7!}{2!\cdot 2!}\]
        \end{tcolorbox}
      \end{enumerate}

   \item A child has $12$ blocks, of which $6$ are black, $4$ are red, $1$ is white, and $1$ is blue.  If the child puts the blocks in a line, how many arrangements are possible?

     \begin{tcolorbox}[%
       enhanced, 
       breakable,
       frame hidden,
       overlay broken = {
         (frame.north west) rectangle (frame.south east);},
       ]
       Let's assume that blocks of the same color are indistinguishable.  To start there are $12$ different blocks, giving a total of $12!$ different arrangements.  But the the $6$ black blocks, and  $4$ red.  Thus there are a total of
       \[\frac{12!}{4!\cdot 6!}\]
       arrangements.
     \end{tcolorbox}

   \item In how many ways can $8$ people be seated in a row if
     \begin{enumerate}[label=\alph*)]
\item there are no restrictions on the seating arrangmeent?

  \begin{tcolorbox}[%
    enhanced, 
    breakable,
    frame hidden,
    overlay broken = {
      (frame.north west) rectangle (frame.south east);},
    ]
    $8!$ different arrangements.
  \end{tcolorbox}

  \item persons $A$ and $B$ must sit next to each other?

    \begin{tcolorbox}[%
      enhanced, 
      breakable,
      frame hidden,
      overlay broken = {
        (frame.north west) rectangle (frame.south east);},
      ]
      Then when ordering the people, we can treat $A$ and $B$ as a group, thus there are $7! \cdot 2!$ different arrangements ($2!$ different ways to organize $A$ and $B$ within their group).
    \end{tcolorbox}

    \item there are $4$ men and $4$ men, and no $2$ men or $2$ womaen can sit next to each other?

      \begin{tcolorbox}[%
        enhanced, 
        breakable,
        frame hidden,
        overlay broken = {
          (frame.north west) rectangle (frame.south east);},
        ]
        This is basically saying that there are two possible arrangements \texttt{MFMFMFMF} \texttt{FMFMFMFM}.  Within each option there are $4!$ ways to arrange the men, and $4!$ ways to arrange the women.  Thus there are a total of $2 \cdot 4! \cdot 4!$ possible different arrangements.  
      \end{tcolorbox}

      \item there are $5$ men, $3$ women and all the men must sit next to eachother?

        \begin{tcolorbox}[%
          enhanced, 
          breakable,
          frame hidden,
          overlay broken = {
            (frame.north west) rectangle (frame.south east);},
          ]
          There are
          \begin{itemize}
          \item $5!$ different ways to arrange the men
          \item $3!$ different ways to arrange the women
          \item After placing the $3$ women, there are $4$ different places to place the group of men.
          \end{itemize}
          Thus there are a total of $5! \cdot 3! \cdot 4$ different possible arrangements.
        \end{tcolorbox}

      \item There are $4$ married couples and each couple must sit together?

        \begin{tcolorbox}[%
          enhanced, 
          breakable,
          frame hidden,
          overlay broken = {
            (frame.north west) rectangle (frame.south east);},
          ]
          There are $4$ different groups, and each group can be arranged $2! = 2$ different ways.  Thus there are a total of
          \[4! \cdot 2^4\]
          different arrangements.
        \end{tcolorbox}
      \end{enumerate}

    \item In how many ways can $3$ novels, $2$ math books, and $1$ chemistry book be arranged on a bookshelf if
      \begin{enumerate}[label=\alph*)]
\item the books can be arranged in any order?

  \begin{tcolorbox}[%
    enhanced, 
    breakable,
    frame hidden,
    overlay broken = {
      (frame.north west) rectangle (frame.south east);},
    ]
    There are $6$ books total, and since order doesn't matter, there are $6!$ ways to arrange them.
  \end{tcolorbox}

  \item the math books must be together and the novels must be together?

    \begin{tcolorbox}[%
      enhanced, 
      breakable,
      frame hidden,
      overlay broken = {
        (frame.north west) rectangle (frame.south east);},
      ]
      Thinking about this in components we have:
      \begin{itemize}
      \item $3$ groups: $3!$ ways to arrange each group
      \item $3!$ ways to arrange the novels
      \item $2!$ ways to arrange the math books
      \item $1!$ ways to arrange chemistry book
      \end{itemize}
      Thus there is a total of
      \[3! \cdot 3! \cdot 2!\]
      ways to arrange the books.
    \end{tcolorbox}

    \item the novels must be together, but the other books can be arranged in any order?

      \begin{tcolorbox}[%
        enhanced, 
        breakable,
        frame hidden,
        overlay broken = {
          (frame.north west) rectangle (frame.south east);},
        ]
        There are 4 blocks: the novels, 2 math books (one block each), and the chemistry book.  But there are also the $3!$ ways of arranging the novels, so a total of
        \[4! \cdot 3!\]
        arrangements.
      \end{tcolorbox}
    \end{enumerate}

  \item Five seperate awards (best scholarship, best leadership qualities, and so on) are to be presented to selected students from a class of 30.  How many different outcomes are possible if
    \begin{enumerate}[label=\alph*)]
\item a student can receive any number of awards?
  \begin{tcolorbox}[%
    enhanced, 
    breakable,
    frame hidden,
    overlay broken = {
      (frame.north west) rectangle (frame.south east);},
    ]
    At each point in choosing an award winner, there is a pool of $30$ students to choose from.
    Thus there is a total of
    \[30^5\]
    different ways to award the awards.
  \end{tcolorbox}

  \item A student can only recieve one award.

    \begin{tcolorbox}[%
      enhanced, 
      breakable,
      frame hidden,
      overlay broken = {
        (frame.north west) rectangle (frame.south east);},
      ]
      Then, at each point of recieving an award, that student is removed from selection pool.  Thus there are a total of
      \[\frac{30!}{25!}\]
      different ways to hand out the awards.
    \end{tcolorbox}
  \end{enumerate}
 \item Consider a group of 20 people. If everyone shakes hands with everyone else, how many handshakes take place?

   \begin{tcolorbox}[%
     enhanced, 
     breakable,
     frame hidden,
     overlay broken = {
       (frame.north west) rectangle (frame.south east);},
     ]
     Each person shakes $19$ hands, and there are a total of $20$ people.  But each handshake is counted twice, so a total of
     \[\frac{20\cdot 19}{2}\]
    total handshakes.
  \end{tcolorbox}
  \item How many $5$ card poker hands are there?

    \begin{tcolorbox}[%
      enhanced, 
      breakable,
      frame hidden,
      overlay broken = {
        (frame.north west) rectangle (frame.south east);},
      ]
      A standard deck of cards has $52$ cards.  This means that there are a total of
      \[\frac{52!}{47!}\]
      ways to choose $5$ different cards.  But the order of the cards in your hand does not matter,
      so we need to get rid of duplicates.  This means that there is a total of
      \[\frac{52!}{47!5!}\]
      legal $5$ card poker hands.
    \end{tcolorbox}

    \item A dance class consists of 22 students, of which 10 are women and $12$ are men.  If $5$ men and $5$ women are to be chosen and then paired off, how many results are possible?

      \begin{tcolorbox}[%
        enhanced, 
        breakable,
        frame hidden,
        overlay broken = {
          (frame.north west) rectangle (frame.south east);},
        ]
        There are
        \begin{itemize}
        \item $12 \choose 5$ ways to choose the men
        \item $10 \choose 5$ ways to choose the women
        \item $5!$ ways different ways of pairing up the $5$ men with $5$ women.  To see this, you can
          imagine the first man has $5$ options, the next has $4$ and so forth.
        \end{itemize}
        Thus there are a total of
        \[\binom{12}{5} \cdot \binom{10}{5} \cdot 5!\]
        different possible results.
      \end{tcolorbox}
\end{enumerate}


\end{document}




